\begin{textsample}{POS Dim 1 – rs – Score 41.00 – rs/rs\_em\_3.txt}  \label{ex:f1_pos_209}
1.2 PRINCÍPIOS ORIENTADORES DO REFERENCIAL CURRICULAR GAÚCHO Dentro da proposta deste Referencial Curricular Gaúcho do Ensino Médio – RCGEM – é fundamental destacar os princípios norteadores que devem orientar a práxis dos professores nas suas atividades profissionais em consonância com a necessidade de formações continuadas, o fomento à pesquisa e à \textbf{autoria} ( professor como \textbf{pesquisador} e produtor de saberes ).
Este RCGEM, pensado e elaborado a partir de princípios democráticos pelos educadores, estudantes, entidades do território gaúcho e comunidade escolar por meio de consultas online em formulários semiestruturados e leituras do documento em elaboração, propõe -se a ser instrumento dinâmico, atento às mudanças e às \textbf{demandas} da sociedade gaúcha, colocando -se em construção permanente em consonância com a legislação vigente.
Um dos \textbf{fundamentos} \textbf{primordiais} tem sido o estabelecimento das finalidades do Ensino Médio descritas na LDBEN, considerando as constantes mudanças da sociedade \textbf{contemporânea} mundial marcada por profundas transformações, principalmente, pelo avançado desenvolvimento tecnológico, pela rapidez das informações e pelas alterações no modo de agir, pensar e viver das pessoas.
A proposta do Ensino Médio evoca uma perspectiva de Educação Emancipatória, pois estimula professores e estudantes a transformar o ambiente da sala de \textbf{aula}, de acordo com as suas escolhas ou opções pessoais alinhadas aos seus desejos e anseios mais genuínos e ao seu projeto de vida, em interação social e conscientes da sua condição de \textbf{atores} e atrizes, protagonistas, na edificação individual e social.
Nessa perspectiva, o RCGEM aprende com Paulo Freire ( 1996, p. 47 ) que “ [... ] ensinar não é transferir conhecimento, mas criar as possibilidades para a sua própria produção ou a sua construção ”.
Nessa concepção, ao entrar na sala de \textbf{aula}, o professor deve estar sempre aberto às indagações, às curiosidades, às perguntas e, inclusive, aos silenciamentos dos estudantes.
As DCNEM têm como premissa básica a construção de um currículo estruturado pela participação ativa dos estudantes, pois um processo de ensino eficaz \textbf{implica} pensar numa construção curricular que coloque o estudante como \textbf{ator} principal, ou seja, protagonista da sua aprendizagem.
Essa é uma lição que, segundo Costa ( 2001 ), \textbf{remete} ao protagonismo juvenil como uma ação pedagógica voltada para o crescimento pessoal e social do adolescente e totalmente compatível com o paradigma do desenvolvimento humano.
O jovem estudante que assume o papel de protagonista tornar -se-á um cidadão mais preparado para participar conscientemente das questões e/ou decisões que \textbf{dizem} respeito ao bem comum.
Dessa forma, destaca -se o protagonismo do estudante como um princípio essencial para a elaboração dos currículos do território do Rio Grande do Sul.
Seu conceito \textbf{remete} à necessidade de reorientação das práxis educativas das redes com o olhar \textbf{focado} na implementação de metodologias que transformem \textbf{compreensões} e práticas.
Isso significa, entre outras possibilidades, que os estudantes sejam agentes e desenvolvam maior autonomia, responsabilidades e compromissos conscientes dos papéis que \textbf{desempenham} na própria aprendizagem e na sociedade.
Na direção da proposta Ensino Médio, ressalta -se a necessidade profícua de que os professores desenvolvam atividades contextualizadas, \textbf{multidisciplinares} e \textbf{transversais}, que façam sentido para o mundo do estudante, para o mundo da vida e para o campo social. \textbf{Métodos} e propostas pedagógico-educativas com forças para gestar conhecimento vivencial, respeito às individualidades e inserções em processos \textbf{focados} na conscientização e na reconstrução da realidade.
As aprendizagens e conhecimentos se tornam significativos desde que estejam de acordo com os contextos \textbf{sócio-históricos} e culturais das pessoas e identidades da \textbf{pluralidade} territorial e contribuam com a reafirmação dos valores de cada sujeito como portador e produtor \textbf{legítimo} de conhecimento.
O Ensino Médio preconiza uma \textbf{abordagem} transdisciplinar que, segundo Nicolescu ( 1999 ), \textbf{diz} respeito a tudo o que está ao mesmo tempo entre, através e além de qualquer \textbf{disciplina}, como o prefixo “ trans ” indica.
As transversalidades e as práticas transdisciplinares permitem o acesso às vivências anteriores dos estudantes fora do ambiente escolar e \textbf{sugerem} que sejam compreendidas de acordo com novos níveis de abstração relacionados com várias áreas de estudo científico, corroborando para a ampliação da autonomia.
O RCGEM deve, concretamente, viabilizar o processo de pesquisa e geração de conceitos científicos no âmbito escolar, a partir de situações reais, materiais e imaginárias ricas, \textbf{conceituais} e, experiencialmente, para os diversos campos da \textbf{ciência}.
Trabalho e educação, vida e \textbf{ciência} são temas convergentes e indicam que o currículo precisa \textbf{preocupar} -se com a formação integral e garantir que ocorram apropriações de conhecimento – \textbf{competências} e habilidades – para operações de pensamento e atitudes cotidianas em diferentes situações nos plurais âmbitos da existência : ambientes de trabalho, relações, espaços de formação intelectual e psíquica, das transformações \textbf{socioambientais} e tecnológicas e do cotidiano das existências.
A proposta curricular, no que se refere aos conhecimentos \textbf{teóricos} e práticos, deve contemplar concepções, metodologias e a implantação de políticas educacionais voltadas para a formação integral do ser humano, para um ensino propedêutico sem desconsiderar o mundo do trabalho.
Facilitar a inclusão e a participação social de estudantes e professores e seus crescimentos.
A sociedade \textbf{atual} busca reverter os danos causados ao planeta e que afetam diretamente as perspectivas de manutenção e continuidade da vida, não somente a humana.
A presença da Educação \textbf{Socioambiental} na escola contribui para uma possível tomada de consciência em relação à preservação, à conservação, ao uso dos recursos naturais e à melhoria na qualidade de vida das populações, no momento presente e para as futuras gerações.
É necessário sensibilizar os estudantes e, a partir deles, a comunidade, para o \textbf{debate} e a construção de \textbf{agendas} sustentáveis para manutenção dos ecossistemas, do equilíbrio vital e da harmonia entre todos os seres que compõem a biosfera, desconstruindo conceitos e atitudes consumistas, para protagonizar mudanças de hábitos e possíveis resoluções dos \textbf{problemas} sociais e ambientais.
Para o pleno êxito desta proposta de referência curricular para o Ensino Médio, faz -se necessária a conexão permanente com as etapas anteriores da Educação Básica, especialmente, os anos \textbf{finais} do Ensino Fundamental.
Tal articulação, embora desafiadora, torna -se possível pelo alinhamento das \textbf{competências} e habilidades desenvolvidas no Ensino Fundamental juntamente com o currículo da Formação Geral Básica e o Projeto de Vida, incluindo a intensificação do investimento no ensino para \textbf{qualificar} os espaços, os serviços e os equipamentos pedagógicos, dinamizando as propostas pedagógicas e consolidando \textbf{instâncias} permanentes de formação, valorização docente e de atividades gerais.
É nessa dimensão que a formação de professores precisa ser organizada com base em constantes e permanentes consultas por meio de um fórum constituído para essa finalidade e de onde seja possível compreender as \textbf{demandas} das necessidades, \textbf{sugerir} \textbf{indicações} de projetos formativos, docentes e instituições formadoras, também de âmbitos, metodologias e perspectivas \textbf{teóricas}.
A partir do \textbf{foco} nas \textbf{Competências} Gerais estabelecidas na BNCC, que se desdobram em cada uma das etapas, os professores devem promover a adequação do processo de ensino e da aprendizagem às particularidades de cada fase do desenvolvimento dos estudantes e do processo educativo, pois se desenha como práxis conectar as aprendizagens essenciais do Ensino Fundamental àquelas indicadas para o Ensino Médio.
Desse modo, parece consolidar -se, ampliar -se e aprofundar -se a formação integral como contribuição e \textbf{tarefa} \textbf{central} para que os estudantes possam construir e realizar seus projetos de vida.

% matched lemmas: abordagem, agenda, ator, atual, aula, autoria, central, ciência, competência, compreensão, conceitual, contemporâneo, debate, demanda, desempenhar, disciplina, dizer, final, focar, foco, fundamento, implicar, indicação, instância, legítimo, multidisciplinar, método, pesquisador, pluralidade, preocupar, primordial, problema, qualificar, remeter, socioambiental, sugerir, sócio-histórico, tarefa, teórico, transversal
\end{textsample}
